% \documentclass{article}
% \usepackage{graphicx} % Required for inserting images

% \title{Rubric for Homework 2}

% \begin{document}

% \maketitle

% \section{$Q_1$ Game Theory}
% \begin{itemize}
%     \item 1
% \end{itemize}


% \end{document}

\documentclass[11pt, letterpaper]{article}

% PACKAGES
\usepackage[margin=1in]{geometry}
\usepackage{amsmath}
\usepackage{amssymb}
\usepackage{graphicx}
\usepackage[utf8]{inputenc}
\usepackage{hyperref}

% DOCUMENT LAYOUT
\setlength{\parindent}{0pt}
\setlength{\parskip}{1em}
\renewcommand{\labelitemi}{\textbullet}
\renewcommand{\labelitemii}{$\circ$}

\hypersetup{
    colorlinks=true,
    linkcolor=blue,
    filecolor=magenta,      
    urlcolor=cyan,
}

% HELPER COMMANDS
\newcommand{\question}[2]{
    \subsection*{Question #1: #2}
}
\newcommand{\points}[1]{
    \textbf{(#1 Points)}
}

\begin{document}

% --- TITLE ---
\begin{center}
    {\Huge \textbf{Midterm Exam Grading Rubric}} \\
    \vspace{0.2cm}
    {\large IEDA 1250: Optimization for Decision Making (Summer, 2026)} \\
    \vspace{0.5cm}

    \hrule
\end{center}

% --- QUESTION 1 ---


% --- QUESTION 1 ---
\question{1}{Linear Programming -- Graphical Method \points{10}}

\begin{itemize}
    \item Correctly identifies that the feasible region is a diamond bounded by
    $0 \leq x_1+2x_2 \leq 6$ and $-2 \leq 2x_1-x_2 \leq 4$. \points{2}

    \item Correctly finds the four vertices of the feasible region:
    \[
    \left(-\frac{4}{5},\frac{2}{5}\right),
    \left(\frac{8}{5},-\frac{4}{5}\right),
    \left(\frac{2}{5},\frac{14}{5}\right),
    \left(\frac{14}{5},\frac{8}{5}\right).
    \]
    \points{4}

    \item Correctly identifies the optimal solution and optimal objective value:
    \[
    x_1^*=\frac{14}{5},\qquad x_2^*=\frac{8}{5},\qquad z^*=\frac{18}{5}.
    \]
    \points{4}
\end{itemize}


% --- QUESTION 2 ---
\question{2}{Linear Programming -- Complementary Slackness \points{20}}

\begin{itemize}
    \item Correctly writes or uses the dual constraints corresponding to the primal LP,
    taking into account that $x_3,x_4,x_5$ are unrestricted in sign. \points{4}

    \item Correctly applies complementary slackness to the primal constraints and obtains
    \[
    2x_1+x_2+x_3+x_4=3.
    \]
    \points{4}

    \item Correctly applies complementary slackness to the dual constraints for $x_1$
    and $x_2$, and concludes that
    \[
    x_1=0,\qquad x_2=0.
    \]
    \points{5}

    \item Correctly substitutes $x_1=0$ and $x_2=0$ into the primal constraints and forms
    the reduced system:
    \[
    \begin{cases}
    x_3+x_4=3,\\
    2x_3+x_5=7,\\
    x_3-x_4+2x_5=7.
    \end{cases}
    \]
    \points{4}

    \item Correctly solves the reduced system and obtains the optimal primal solution:
    \[
    x_1^*=0,\quad x_2^*=0,\quad x_3^*=2,\quad x_4^*=1,\quad x_5^*=3.
    \]

    the optimal objective value:
    \[
    z^*=6.
    \]
    
    \points{3}

    % \item Correctly reports the optimal objective value:
    % \[
    % z^*=6.
    % \]
    % \points{1}
\end{itemize}



\end{document}